\scanpage{62}
Etude numérique des catégories tannakiennes.

a) Soit $\handwritten{\mathcal{M}}$ une catégorie abélienne $k$-linéaire avec $\operatorname{Hom}$ de dimension finie ($k$ un corps). Si $M$ est un objet simple de $\handwritten{\mathcal{M}}$, la sous-catégorie abélienne de $\handwritten{\mathcal{M}}$ qu'il engendre est équivalente à la catégorie des vectoriels de dimension finie sur le corps $K=\operatorname{End}(M)$. On voit ain\unsure{\deleted{is}si|\deleted{si}is} que l'\unsure{\deleted{i}o|\deleted{o}i}bjet $M$ reste semi-simple par une extension $k'/k$ sss $K\otimes_k k'$ est un comp\unsure{\deleted{l}o|\deleted{o}l}sé \deleted{de corps} d'algèbres de matrices (i.e. est une algèbre semi-simple). Ceci est v\deleted{rai}\handwritten{érifié} pour tout $k'$ sss c'est vrai pour \deleted{$k'$ \unsure{c}lôtur\unsure{e}}\inserted{\handwritten{\uline{une} ext.}\begin{tikzpicture}[x=1cm,y=1cm,baseline=-.1cm]\draw[thin] (0,0)--(.8,0).. controls (.87,0) and (.88,.1)..(.91,.12)--(1.01,-.08);\end{tikzpicture}} \uline{parfaite} de $k$, \deleted{i}\unsure{\deleted{e}o|\deleted{o}e}u sss le centre \deleted{de} $Z$ de $K$ est séparable sur $k$\deleted{.}\begin{tikzpicture}[x=1cm,y=1cm,baseline=-.1cm]\draw[thin] (-.1,-.25).. controls (0,0) and (.15,.1)..(.25,.35);\end{tikzpicture} On dit alors que $M$ est absolument semi-simple (à généraliser au cas n'est pas supposé simple).

Le groupe $K(\handwritten{\mathcal{M}})$ est le groupe libre engendré par l'ensemble $\handwritten{\Sigma}(\handwritten{\mathcal{M}})$ des classes d'objets simples de $\handwritten{\mathcal{M}}$. Si $k'$ est une extension de $k$, la fa\unsure{\deleted{v}ç|\deleted{ç}v}\deleted{\unsure{e|c}}on dont un objet simple $M$ de $\handwritten{\mathcal{M}}$ tel que $M_{k'}$ soit semi-simple se décompose enn objets simples se voit sur la structure de $K\otimes_k k'=K'$ : si $K'$ est le produit d'algèbres $M_{n_i}(K'_i)$, où les $K'_i$ sont des corps gauches,\inserted{$1\leqslant i\leqslant r$,\begin{tikzpicture}[x=1cm,y=1cm,baseline=-.1cm]\draw[thin] (0,.15).. controls (.1,-.1) and (-.2,-.1)..(-1,-.1)--(-2.1,-.1).. controls (-2.3,-.1) and (-2.35,-.2)..(-2.35,-.5);\end{tikzpicture}} alors $M_{k'}$ se décompose en $r$ composantsisotypiques (correspondants aux $K'_i$) chacun ayant $n_i$ composants simples\inserted{$M'_i$\begin{tikzpicture}[x=1cm,y=1cm,baseline=-.1cm]\draw[thin] (-.1,-.25).. controls (0,0) and (.15,.1)..(.25,.35);\end{tikzpicture}}. Si $k'$ est une extension quasi-galoisienne de $k$, les classes des $M'_i$ sont conjuguées entre elles par l'action de $\operatorname{Gal}(k'/k)=\handwritten{\pi}$. On trouve alors une bijection canonique (si l'hypothèse faite est vérifiée pour tout objet simple \deleted{de} $M$ de $\handwritten{\mathcal{M}}$, p.ex. si $k'/k$ est même galoisienne)
\[
\handwritten{\Sigma(\mathcal{M})\simeq\Sigma(\mathcal{M}')/\pi}\ .
\]
Il y aurait lieu de généraliser au cas où on ne suppose pas nécessairement que les $K\otimes_k k'$ sont semi-simples ....

\inserted{Si $\handwritten{\mathcal{M}}$ est semi-simple,\begin{tikzpicture}[x=1cm,y=1cm,baseline=-.1cm]\draw[thin] (0,.15).. controls (.1,-.1) and (-.2,-.1)..(-1,-.1)--(-2.1,-.1).. controls (-2.3,-.1) and (-2.35,-.2)..(-2.35,-.5);\end{tikzpicture}} \deleted{L}\handwritten{l}a catégorie $\handwritten{\mathcal{M}}$ est connue à équivalence près quand on connait l'ensemble $\handwritten{\Sigma}=\handwritten{\Sigma(\mathcal{M})}$, et l'application
\[
\handwritten{\sigma\longmapsto (Z_\sigma,\xi_\sigma\in}\operatorname{Br}(\unsure{\deleted{k)}|\deleted{k}}\handwritten{Z_\sigma})\handwritten{)}
\]
associant à toute classe\inserted{\handwritten{$\sigma$\begin{tikzpicture}[x=1cm,y=1cm,baseline=-.1cm]\draw[thin] (-.1,-.25).. controls (0,0) and (.15,.1)..(.25,.35);\end{tikzpicture}}} d'objet simple $\unsure{\deleted{S}M|\deleted{M}S}$ de $S$ l\unsure{\deleted{a}e|\deleted{e}a} \deleted{classe de}\inserted{\handwritten{centre $Z_\sigma$}} $\operatorname{End}(\deleted{\unsure{S}}\handwritten{\mathcal{M}})$
% archive sheet number: 61
